% !TEX root = ../../../main.tex
% 所属：第四部分：Fourier 分析 / 第3章：Fourier级数的一些应用
% 节：3.3 构造一个连续但是处处不可微的函数
% 正文按原文件顺序保存；由 ../chapter3.tex 调用。
\section{构造一个连续但是处处不可微的函数}

这里,我们对 Weierstrass 的一般构造取一个特例,并且用另外方式证明它.

 我们需要先了解一些相关的知识:部分和,Cesàro 和,推延平均.

\begin{definition}
定义\textbf{推延平均}
\[
\Delta_N(g)(x)=2\sigma_{2N}(g)(x)-\sigma_N(g)(x).
\]
\end{definition}

由 Fejér 核 $F_N$ 的卷积表示
\[
\sigma_N(g)=g*F_N,
\]
立即可得
\[ \Delta_N(g)=g*[2F_{2N}-F_N]. \]


若 $g(x)\sim\sum a_ne^{inx}$,则利用上述定义直接计算可得
\begin{enumerate}
\item $S_N$ 在 $|n|\leq N$ 时,对 $a_ne^{inx}$ 贡献 1 倍,在 $|n|>N$ 时,贡献 0 倍;
\item $\sigma_N$ 在 $|n|\leq N$ 时,对 $a_ne^{inx}$ 贡献 $1-\dfrac{|n|}{N}$ 倍,在 $|n|>N$ 时,贡献 0 倍;
\item $\Delta_N$ 在 $|n|\leq N$ 时,对 $a_ne^{inx}$ 贡献 1 倍,在 $N<|n|\leq 2N$ 时,贡献 $2(1-\dfrac{|n|}{2N})$ 倍,在 $|n|>2N$ 时,贡献 0 倍.
\end{enumerate}
\begin{center}
	\includegraphics[width=0.7\linewidth]{"figure/P4/fig for 3.3"}
\end{center}
\begin{theorem}

若 $0<\alpha<1$,则函数
\[
f(x)=f_\alpha(x)=\sum_{n=0}^{\infty}2^{-n\alpha}e^{i2^n x}
\]
连续但是处处不可微.
\end{theorem}

由级数的绝对收敛性,连续性是显然的. $f$ 的关键性质是它的大多数 Fourier 系数为零,像上式那样缺少很多项的 Fourier 级数称为一个\textbf{缺项 Fourier 级数}.

 对于像 $f$ 这样的缺项级数,推延平均与部分和是基本相等的. 特别地,对函数 $f=f_\alpha$ 有

\[ S_N(f)=\Delta_{N'}(f), \]

其中 $N'$ 是不大于 $N$ 的形如 $2^k$ 的最大整数. 观察 $f$ 的定义,这一点是显然的.

下面使用反证法证明,即假设对于某些 $x_0$,$f'(x_0)$ 存在. 先建立一个关于 Cesàro 平均导数的估计.

\begin{lemma}
设 $g$ 为任意在 $x_0$ 处可微的连续函数,则其 Cesàro 平均满足
\[
\sigma_N(g)'(x_0)=O(\log N).
\]
因此 $\Delta_N(g)'(x_0)=O(\log N)$.
\label{Cesàro平均导数的对数增长}
\end{lemma}

\begin{proof}
由卷积公式 $\sigma_N(g)=g*F_N$ 与 $g$ 在 $x_0$ 处的可微性,有
\[
\sigma_N(g)'(x_0)
=\frac1{2\pi}\int_{-\pi}^{\pi}F_N'(t)g(x_0-t)\dd t.
\]
又因为 $F_N$ 的周期性给出 $\displaystyle\int_{-\pi}^{\pi}F_N'(t)\dd t=0$,故可减去常数 $g(x_0)$:
\[
\sigma_N(g)'(x_0)
=\frac1{2\pi}\int_{-\pi}^{\pi}F_N'(t)\bigl[g(x_0-t)-g(x_0)\bigr]\dd t.
\]

假设 $g$ 在 $x_0$ 处可微,则对 $t$ 充分小有 $\abs{g(x_0-t)-g(x_0)}\leq C\abs{t}$. 因此
\[
\abs{\sigma_N(g)'(x_0)}
\leq\frac C{2\pi}\int_{-\pi}^{\pi}\abs{F_N'(t)}\abs{t}\dd t.
\]

现在对 $\abs{F_N'(t)}$ 建立两个估计.
\begin{itemize}
\item 由 Fejér 核的 Fourier 形式 $F_N(t)=\sum\limits_{|k|\leq N}(1-\dfrac{\abs{k}}{N})e^{ikt}$ 求导得
\[
\abs{F_N'(t)}\leq\sum_{|k|\leq N}\abs{k}(1-\dfrac{\abs{k}}{N})\leq AN^2.
\]
\item 利用 $\sin(\dfrac{Nt}{2})\leq CN\abs{t}$ 与 $\sin(\dfrac t2)\geq c\abs{t}$ (对 $\abs{t}\leq\pi$),可以建立
\[
\abs{F_N'(t)}\leq\frac A{\abs{t}^2}.
\]
\end{itemize}

综合两个估计并对积分区间拆分:
\[
\begin{aligned}
\abs{\sigma_N(g)'(x_0)}
&\leq\frac C{2\pi}\left[\int_{\abs{t}\geq1/N}\abs{F_N'(t)}\abs{t}\dd t
+\int_{\abs{t}\leq1/N}\abs{F_N'(t)}\abs{t}\dd t\right]\\
&\leq\frac{CA}{2\pi}\left[\int_{\abs{t}\geq1/N}\frac{\dd t}{\abs{t}}
+N^2\int_{\abs{t}\leq1/N}\abs{t}\dd t\right]\\
&=\frac{CA}{2\pi}\bigl[O(\log N)+O(1)\bigr]=O(\log N).
\end{aligned}
\]

由 $\Delta_N(g)=2\sigma_{2N}(g)-\sigma_N(g)$ 立即得 $\Delta_N(g)'(x_0)=O(\log N)$.
\end{proof}

\begin{lemma}
若 $2N=2^k$ (即 $N=2^{k-1}$),则
\[
\Delta_{2N}(f)(x)-\Delta_N(f)(x)=2^{-k\alpha}e^{i2^k x}.
\]
\label{缺项Fourier级数的差分}
\end{lemma}

\begin{proof}
考虑$f$的定义 ，这一点显然
\end{proof}

\begin{proof}[定理的证明]
连续性由绝对收敛性即得. 只需证 $f$ 在圆周上处处不可微.

设存在 $x_0$,使 $f'(x_0)$ 存在. 则 $f$ 在 $x_0$ 处可微,且 $f$ 显然是连续函数. 由引理 \ref{Cesàro平均导数的对数增长},
\[
\abs{[\Delta_{2N}(f)-\Delta_N(f)]'(x_0)}=O(\log N).
\]

另一方面,任取 $k\geq1$,令 $N=2^{k-1}$,从而 $2N=2^k$. 由引理 \ref{缺项Fourier级数的差分},
\[
\Delta_{2N}(f)(x)-\Delta_N(f)(x)=2^{-k\alpha}e^{i2^k x}.
\]
对 $x$ 求导得
\[
[\Delta_{2N}(f)-\Delta_N(f)]'(x)=i2^k\cdot2^{-k\alpha}e^{i2^k x}.
\]
特别地,
\[
\abs{[\Delta_{2N}(f)-\Delta_N(f)]'(x_0)}=2^k\cdot2^{-k\alpha}=2^{k(1-\alpha)}=2N^{1-\alpha}.
\]

由于 $0<\alpha<1$,故 $1-\alpha>0$,从而 $2N^{1-\alpha}\to\infty$,其增长速度比 $\log N$ 快得多. 这与引理 \ref{Cesàro平均导数的对数增长} 矛盾. 因此 $f'(x_0)$ 不存在,即 $f$ 处处不可微.
\end{proof}
